% !TEX root = ./rapport.tex
% ══════════════════════════════════════════════════════════════════════
%  ÉTUDE THÉORIQUE — Prédiction de l'attrition
%  Inclus dans chapitre 5 avant le Sprint 7
% ══════════════════════════════════════════════════════════════════════

\subsection{Étude théorique~: Prédiction de l'attrition du personnel}

Lors de l'analyse des besoins du sprint~7, la conception d'un module
capable d'anticiper les départs volontaires s'est imposée~: l'objectif n'était pas de
prédire avec certitude qui allait partir, mais d'identifier les profils présentant
des signaux comportementaux qui, dans la littérature RH, sont associés à une
probabilité de départ plus élevée. Ce module devait être opérationnel dès le
déploiement initial de Synapse, sans historique de départs numérisé.

\subsection{Qu'est-ce que l'attrition~?}

L'attrition désigne le départ non remplacé d'employés d'une organisation. Elle
représente un coût significatif~: au-delà du recrutement, la perte de compétences
accumulées et la baisse de productivité durant la période d'intégration du
remplaçant constituent un coût souvent sous-estimé. Détecter les signaux
précurseurs d'un départ permet d'agir en amont~: entretien de carrière,
revalorisation, ou affectation à un projet plus stimulant.

\subsection{Choix de l'approche~: scoring par règles}

Deux grandes familles d'approches existent pour modéliser l'attrition.

\begin{table}[H]
\centering
\caption{Comparaison des approches de modélisation de l'attrition}
\label{tab:attrition_approaches_theory}
\begin{tabularx}{\textwidth}{|l|L|L|}
\hline
\rowcolor{headerblue!55}
\textbf{Approche} & \textbf{Conditions requises} & \textbf{Application à Synapse} \\
\hline
\rowcolor{rowgray}
Scoring par règles métier & Indicateurs observables, règles connues & \textbf{Retenue}~: applicable dès le déploiement \\
\hline
Machine Learning supervisé & Historique de départs étiquetés (200+ cas minimum) & Perspective~: à envisager après 12--24 mois d'exploitation \\
\hline
\rowcolor{rowgray}
Modèles de survie & Données de durée d'ancienneté complètes & Avancée, non prioritaire \\
\hline
\end{tabularx}
\end{table}

L'approche par machine learning supervisé a été écartée car Arabsoft ne dispose pas, au
moment du déploiement, d'un historique suffisant de départs numérisés pour entraîner
un modèle fiable. Un modèle entraîné sur quelques dizaines d'exemples produirait
des prédictions peu exploitables. L'approche par règles, fondée sur des indicateurs
directement mesurables dans la base Oracle de Synapse, est donc la plus cohérente
avec la réalité du terrain.

\subsection{Facteurs retenus et leur signification}

Les quatre facteurs retenus sont documentés dans la littérature comme des signaux précurseurs de départ\textsuperscript{[26,27]}~: l'\textbf{ancienneté courte} (faible attachement organisationnel)~; un \textbf{volume d'absences élevé} (corrélé à l'insatisfaction)~; un \textbf{taux de rejet élevé des demandes} (perception d'injustice diminuant l'engagement)~; et l'\textbf{absence de projets actifs} (marginalisation professionnelle). Ces quatre indicateurs sont calculables directement depuis les tables Oracle de Synapse (\texttt{EMPLOYEES}, \texttt{LEAVE\_REQUESTS}, \texttt{V\_ALL\_DEMANDES}, \texttt{PROJECT\_MEMBERS}).

\subsection{Limites reconnues}

Ce module produit un score de risque, non une prédiction certifiée. Plusieurs facteurs d'attrition reconnus (rémunération, évolutions de
carrière, dynamique managériale) ne sont pas encore intégrés au schéma de Synapse.
De plus, le scoring est déterministe~: deux employés avec les mêmes indicateurs
obtiennent le même score, sans tenir compte des interactions entre facteurs. Ces
limites sont assumées et définissent la feuille de route naturelle du module~: à
mesure que l'historique s'accumule et que le schéma s'enrichit, une transition vers
un modèle supervisé deviendra viable.
